% Computer Science A --- Spring Framework (Core, MVC, Data). Boot-specific items live in \texttt{springboot.tex}.
% Free-response bank at the front is migrated from \texttt{cs-a/06-practice.tex} (Spring core, Web/MVC, Data JPA only).

\runningheader{Computer Science A}{Spring Framework}{Page \thepage\ of \numpages}

\begingradingrange{csaspring}

\uplevel{\par\textbf{Free response: Spring core, Web/MVC, and Data JPA}\par
\textit{These items mirror the unit bank; use the boxed space for student work. The printed key expands each idea for study.}\par}

\uplevel{\par\textbf{Spring core: IoC and beans}\par}

\question[4]
What is Inversion of Control? What are its advantages?
\begin{solutionorbox}[1.45in]
  \textbf{Inversion of Control (IoC)} means the framework owns the ``main'' flow: it constructs objects, resolves collaborators, and calls your code at well-defined extension points, instead of your classes using \texttt{new} everywhere to wire themselves.

  \noindent\textbf{Advantages (typical teaching list):}
  \begin{itemize}
    \item \textbf{Looser coupling:} depend on interfaces or abstractions; the container supplies implementations.
    \item \textbf{Easier testing:} inject fakes, stubs, or mocks without subclassing production types.
    \item \textbf{Centralized composition:} wiring rules live in configuration (annotations, Java config, or XML), so collaboration graphs are easier to read and change.
    \item \textbf{Clearer responsibilities:} domain code focuses on behavior; infrastructure concerns (transactions, HTTP, persistence adapters) attach declaratively.
  \end{itemize}
\end{solutionorbox}

\question[4]
What is Dependency Injection?
\begin{solutionorbox}[1.2in]
  \textbf{Dependency injection (DI)} is a form of IoC where an object receives its collaborators from \emph{outside} the object rather than constructing them internally.

  \noindent\textbf{Spring's role:} the IoC container resolves dependencies by type and/or name and injects them using \textbf{constructor} injection, \textbf{setter} injection, or (less preferred for required deps) \textbf{field} injection. The key idea is \emph{push dependencies in} instead of \emph{pulling with \texttt{new}}.
\end{solutionorbox}

\question[4]
What is the difference between XML-based and annotation-based configuration?
\begin{solutionorbox}[1.35in]
  \textbf{XML-based configuration} declares beans, property values, and wiring in external files (\texttt{<bean>}, \texttt{<context:component-scan>}, etc.). It keeps wiring separate from Java source and can be convenient for ops-owned toggles or legacy integration.

  \textbf{Annotation-based configuration} colocates metadata with code (\texttt{@Component}, \texttt{@Service}, \texttt{@Configuration}, \texttt{@Bean}, \texttt{@Autowired}, \ldots), often discovered via component scanning.

  \noindent\textbf{Modern practice:} Java configuration and annotations dominate new Spring code; XML remains valid and is still used where teams standardize on it or integrate older modules. The same container can mix styles during migration.
\end{solutionorbox}

\question[4]
What is a Bean?
\begin{solutionorbox}[1.05in]
  In Spring, a \textbf{bean} is an object whose lifecycle is managed by the \textbf{IoC container}: construction, dependency wiring, initialization callbacks, and (for singletons) a well-defined scope in the \texttt{ApplicationContext}.

  Beans are identified by \textbf{name} and \textbf{type(s)}; configuration (XML, Java config, or stereotype scanning) tells the container which classes to instantiate and how to inject their collaborators. This is distinct from the older JavaBeans naming convention, though Spring can still work with simple property-based objects.
\end{solutionorbox}

\question[4]
What is the Spring Framework? What are its benefits?
\begin{solutionorbox}[1.35in]
  The \textbf{Spring Framework} is a mature Java ecosystem for building enterprise and web applications. Its center of gravity is the \textbf{IoC container} and \textbf{dependency injection}, extended by modules for web (\textbf{Spring MVC} / WebFlux), data access (\textbf{Spring JDBC}, \textbf{Spring Data}), integration, security, AOP, and more.

  \noindent\textbf{Benefits often cited in courses:}
  \begin{itemize}
    \item \textbf{Portable programming model} across servlet containers, embedded servers, and cloud deployments.
    \item \textbf{Composable modules:} adopt only what you need (web without full JPA, etc.).
    \item \textbf{Declarative cross-cutting:} transactions (\texttt{@Transactional}), validation, security filters, and AOP without scattering boilerplate.
    \item \textbf{Large ecosystem and documentation,} plus first-class test support (\texttt{spring-test}, \texttt{MockMvc}, test slices in Boot).
  \end{itemize}
\end{solutionorbox}

\question[4]
What does the Spring Core module provide?
\begin{solutionorbox}[1.2in]
  \textbf{Spring Core} supplies the foundational IoC container: \texttt{BeanFactory} / \texttt{ApplicationContext}, bean lifecycle and scopes, \textbf{dependency injection}, event publication, internationalization helpers, and the \textbf{resource abstraction} (classpath, file, URL resources).

  It also includes \textbf{SpEL} (Spring Expression Language) for property placeholders and dynamic values. Other Spring modules (MVC, Data, Integration) build on this core.
\end{solutionorbox}

\question[4]
How could you create an IoC Container?
\begin{solutionorbox}[1.2in]
  You obtain a Spring \textbf{\texttt{ApplicationContext}}---the usual face of the IoC container---by one of these patterns:
  \begin{itemize}
    \item \textbf{Spring Boot:} \texttt{SpringApplication.run(MyApp.class, args)} builds the context from your \texttt{@SpringBootApplication} class, auto-configuration, and component scan.
    \item \textbf{Plain Spring (Java config):} \texttt{new AnnotationConfigApplicationContext(AppConfig.class)} loads \texttt{@Configuration} classes and scanned components.
    \item \textbf{XML (classic):} \texttt{new ClassPathXmlApplicationContext("context.xml")} loads bean definitions from XML.
  \end{itemize}
  Once the context starts, beans are available for injection or explicit \texttt{getBean} lookup.
\end{solutionorbox}

\newpage

\uplevel{\par\textbf{Spring Web and MVC}\par}

\question[4]
Describe MVC architecture.
\begin{solutionorbox}[1.35in]
  \textbf{Model--View--Controller (MVC)} separates concerns in a user-facing application:
  \begin{itemize}
    \item \textbf{Model:} domain data and rules (entities, DTOs, aggregates); the ``state'' the user cares about.
    \item \textbf{View:} presentation for humans (HTML templates, Thymeleaf/JSP) or, in many APIs, the serialized representation produced for the client.
    \item \textbf{Controller:} translates HTTP requests into operations on the model and selects a view or response payload; it should stay thin and delegate heavy logic to services.
  \end{itemize}
  In \textbf{Spring MVC}, \texttt{DispatcherServlet} maps requests to controller methods, invokes services/repositories, and uses view resolvers or message converters to build the HTTP response.
\end{solutionorbox}

\question[4]
Describe commonly used Spring Web annotations.
\begin{solutionorbox}[1.55in]
  Spring Web maps HTTP to Java methods with \textbf{mapping} annotations and binds inputs with \textbf{parameter} annotations.

  \noindent\textbf{Mapping (examples):} \texttt{@RequestMapping} (any method unless narrowed), \texttt{@GetMapping}, \texttt{@PostMapping}, \texttt{@PutMapping}, \texttt{@DeleteMapping}, \texttt{@PatchMapping}.

  \noindent\textbf{Stereotypes:} \texttt{@Controller} (classic MVC + optional view names), \texttt{@RestController} (\texttt{@Controller} + default \texttt{@ResponseBody} for APIs).

  \noindent\textbf{Binding inputs:} \texttt{@RequestParam} (query/form), \texttt{@PathVariable} (URI templates), \texttt{@RequestHeader}, \texttt{@CookieValue}, \texttt{@RequestBody} (body to object via HttpMessageConverters).

  \noindent\textbf{Responses and errors:} \texttt{@ResponseBody}, \texttt{ResponseEntity}, \texttt{@ResponseStatus}, \texttt{@ExceptionHandler}, \texttt{@ControllerAdvice} / \texttt{@RestControllerAdvice}.
\end{solutionorbox}

\question[4]
What does the Spring Web module provide?
\begin{solutionorbox}[1.2in]
  \textbf{Spring Web (MVC)} provides \texttt{DispatcherServlet}, handler mappings and adapters, argument resolvers, \textbf{HTTP message conversion} (JSON/XML, etc.), static resource handling, \textbf{CORS} support, locale/theme resolution for classic MVC, and integration with servlet containers.

  For reactive stacks, \textbf{Spring WebFlux} offers a parallel model (\texttt{RouterFunction}, \texttt{WebClient}); the exam bank here focuses on servlet-style MVC unless your course emphasizes both.
\end{solutionorbox}

\question[4]
What are some annotations that can be used to get information from the request?
\begin{solutionorbox}[1.35in]
  Common request-binding annotations include:
  \begin{itemize}
    \item \texttt{@RequestParam} for query string and form fields.
    \item \texttt{@PathVariable} for URI template segments in the mapping pattern (e.g.\ \texttt{/orders/\{id\}}).
    \item \texttt{@RequestHeader} and \texttt{@CookieValue} for metadata.
    \item \texttt{@RequestBody} to deserialize the body (often JSON) to a POJO or \texttt{String}.
    \item \texttt{HttpServletRequest} / \texttt{HttpServletResponse} parameters when you need low-level servlet access (discouraged for routine business code).
  \end{itemize}
\end{solutionorbox}

\question[4]
What can we use to customize the response sent back to the client?
\begin{solutionorbox}[1.25in]
  Controllers customize HTTP responses using:
  \begin{itemize}
    \item \textbf{Return types:} domain objects with \texttt{@ResponseBody} / \texttt{@RestController}, or \texttt{ResponseEntity} to set status, headers, and body together.
    \item \textbf{Status and headers:} \texttt{@ResponseStatus}, \texttt{ResponseEntity\#status}, \texttt{HttpServletResponse}, or \texttt{ServerHttpResponse} in WebFlux.
    \item \textbf{Message converters:} configure Jackson/XML modules for serialization rules.
    \item \textbf{Global error shaping:} \texttt{@ExceptionHandler} and advice classes map exceptions to problem bodies instead of generic HTML error pages.
  \end{itemize}
\end{solutionorbox}

\question[4]
How does the \texttt{@ExceptionHandler} annotation work?
\begin{solutionorbox}[1.3in]
  Methods annotated with \texttt{@ExceptionHandler(ExceptionType.class)} run when a matching exception propagates out of controller (or advice) invocation. The method can build a body, choose an HTTP status, and log context---returning \texttt{ResponseEntity}, a POJO (with \texttt{@RestControllerAdvice}), or a view name (classic MVC).

  \noindent\textbf{Scope:} handlers can live on a single \texttt{@Controller}/\texttt{@RestController}, or centrally in a class annotated \texttt{@ControllerAdvice} / \texttt{@RestControllerAdvice} for reuse across many controllers.
\end{solutionorbox}

\newpage

\uplevel{\par\textbf{Spring Data JPA}\par}

\question[4]
What configuration is needed in the \texttt{.properties} or \texttt{.yaml} file to use the Spring Data module?
\begin{solutionorbox}[1.35in]
  For JPA with Spring Boot auto-configuration, you typically declare a \textbf{datasource} (\texttt{spring.datasource.url}, \texttt{username}, \texttt{password}, \texttt{driver-class-name} or the JDBC URL form that implies the driver) and \textbf{JPA/Hibernate} settings such as \texttt{spring.jpa.hibernate.ddl-auto}, dialect (often inferred), show-sql flags, and naming strategies.

  Spring Boot then creates a \texttt{DataSource}, \texttt{EntityManagerFactory}, and transaction manager beans; your \texttt{@Entity} classes and repository interfaces complete the picture. Exact keys vary slightly between \texttt{properties} and \texttt{yaml}, but the prefixes are the same.
\end{solutionorbox}

\question[4]
What is the JPA API? What does it provide?
\begin{solutionorbox}[1.3in]
  The \textbf{Java Persistence API (JPA)} is the standard ORM API in Java. It defines how relational rows map to objects using annotations or XML descriptors (\texttt{@Entity}, \texttt{@Table}, relationships, \texttt{@Id}, \texttt{@Version}, \ldots).

  \noindent\textbf{Runtime API:} \texttt{EntityManager} performs \texttt{persist}, \texttt{merge}, \texttt{remove}, \texttt{find}, queries (\texttt{TypedQuery}, Criteria API), and flush/clear lifecycle operations. \textbf{Portability} across providers (Hibernate, EclipseLink, etc.) is a major goal, though vendor-specific features exist when you opt in.
\end{solutionorbox}

\question[4]
What is an ORM?
\begin{solutionorbox}[1.05in]
  \textbf{Object--Relational Mapping (ORM)} bridges the \textbf{object graph} in memory and \textbf{relational tables} in the database: classes map to tables, fields to columns, associations to foreign keys and join tables.

  The ORM provider generates SQL for common patterns, tracks changes to managed entities, and handles caching and flushing policies. It reduces repetitive JDBC for CRUD and navigational access, at the cost of needing to understand fetch strategies and SQL that the ORM emits.
\end{solutionorbox}

\question[4]
What is Spring Data JPA?
\begin{solutionorbox}[1.2in]
  \textbf{Spring Data JPA} is a Spring Data module that eliminates most hand-written DAO boilerplate for JPA. You declare a repository \textbf{interface} extending \texttt{JpaRepository<Entity, Id>} (or related), and Spring generates a proxy implementation at runtime.

  It supports \textbf{derived query methods} from method names, explicit \texttt{@Query} (JPQL or native SQL), paging and sorting (\texttt{Pageable}, \texttt{Sort}), specifications, and integration with Spring's declarative \texttt{@Transactional} boundaries.
\end{solutionorbox}

\question[4]
How do you create a repository with Spring Data JPA?
\begin{solutionorbox}[1.2in]
  \begin{enumerate}
    \item Define a persistence \texttt{@Entity} with an \texttt{@Id}.
    \item Declare a \textbf{public interface} in a package Spring scans, e.g.\ \texttt{public interface OrderRepository extends JpaRepository<Order, Long> \{\}}.
    \item Ensure the JPA vendor and datasource are configured; Spring Data's \texttt{@EnableJpaRepositories} (often implied by Boot) picks up the interface and provides the implementation proxy.
  \end{enumerate}
  No manual \texttt{implements} class is required for the basic pattern; custom behavior can be added via default methods, fragments, or \texttt{@Query}.
\end{solutionorbox}

\question[4]
What methods does the \texttt{JpaRepository} interface provide?
\begin{solutionorbox}[1.35in]
  \texttt{JpaRepository<T, ID>} extends \texttt{PagingAndSortingRepository} and \texttt{CrudRepository}, so you inherit CRUD operations such as \texttt{save}, \texttt{saveAll}, \texttt{findById}, \texttt{existsById}, \texttt{findAll}, \texttt{deleteById}, \texttt{delete}, and paging APIs (\texttt{findAll(Pageable)}).

  \texttt{JpaRepository} adds JPA-oriented helpers such as \texttt{flush}, \texttt{saveAndFlush}, \texttt{deleteInBatch}/\texttt{deleteAllInBatch} (names evolved across Spring Data versions---consult your version's Javadoc). The teaching point is: \emph{most common persistence operations are already on the interface}.
\end{solutionorbox}

\question[4]
What are custom queries and give me an example of how you would write a custom method.
\begin{solutionorbox}[1.45in]
  A \textbf{custom query} is any repository method whose SQL/JPQL is not fully implied by a simple derived name. You supply the query explicitly or implement logic in a custom fragment.

  \noindent\textbf{Example (JPQL with \texttt{@Query}):}
  \begin{lstlisting}[style=java]
@Query("select o from Order o where o.status = :status and o.total > :minTotal")
List<Order> findOpenLarge(@Param("status") OrderStatus status,
                          @Param("minTotal") BigDecimal minTotal);
  \end{lstlisting}

  \noindent Alternatives: native SQL with \texttt{nativeQuery = true}, \texttt{Specification} / Criteria API for dynamic predicates, or a custom interface + \texttt{impl} class wired with \texttt{repositoryBaseClass} patterns for complex SQL.
\end{solutionorbox}

\question[4]
What is the \texttt{@Transactional} annotation? What are its benefits?
\begin{solutionorbox}[1.35in]
  Spring's \texttt{@Transactional} declares a \textbf{transactional boundary}, usually on service-layer methods that call repositories. The framework starts (or joins) a transaction, commits on success, and \textbf{rolls back on failures} according to rollback rules (defaults favor rollback on unchecked exceptions).

  \noindent\textbf{Benefits:}
  \begin{itemize}
    \item Removes hand-written \texttt{try}/\texttt{commit}/\texttt{rollback} JDBC noise.
    \item Centralizes propagation and isolation policies (\texttt{Propagation.REQUIRED}, \texttt{readOnly = true} for queries, etc.).
    \item Keeps resource lifecycle tied to thread-bound \texttt{EntityManager} / JDBC usage in typical Spring JPA setups.
  \end{itemize}
\end{solutionorbox}

\newpage

\question[4]
Spring is
\begin{choices}
  \choice A dependency injection framework for JavaScript.
  \choice A mocking library for Java.
  \choice A JVM based programming language.
  \CorrectChoice A dependency injection framework for Java.
\end{choices}
\begin{solution}
  The \textbf{Spring Framework} for Java provides \textbf{dependency injection}, declarative middleware (transactions, AOP), and modules such as Spring MVC and Spring Data. It is not a language, not JavaScript-focused, and not primarily a mocking library.
\end{solution}

\question[4]
A Spring bean is
\begin{choices}
  \choice The same as a Java bean.
  \CorrectChoice A class which is constructed and managed by Spring.
  \choice Only classes defined in the Spring source code.
  \choice Any bean annotated with \texttt{@GetMapping}.
\end{choices}
\begin{solution}
  A \textbf{bean} is an object whose lifecycle (construction, wiring, destruction) is handled by the Spring \textbf{IoC container}. JavaBeans are a naming/convention pattern; \texttt{@GetMapping} maps HTTP, not bean identity.
\end{solution}

\question[4]
Which of the following is \textbf{not} a way to get a bean from Spring?
\begin{choices}
  \choice Using the \texttt{getBean} method on the \texttt{ApplicationContext} or \texttt{BeanFactory}.
  \choice Using \texttt{@Autowired} with a valid bean definition.
  \choice Using an XML configuration with a valid bean reference.
  \CorrectChoice Using the \texttt{@Aspect} annotation before the declaration field, setter, or constructor.
\end{choices}
\begin{solution}
  Lookup uses the context (\texttt{getBean}), injection (\texttt{@Autowired} / constructor injection), or XML references. \texttt{@Aspect} declares AOP aspects, not bean retrieval. (Source used \texttt{@Autowire}; standard Spring spelling is \texttt{@Autowired}.)
\end{solution}

\question[4]
Which of the following is \textbf{not} a valid stereotype annotation?
\begin{choices}
  \choice \texttt{@Component}
  \CorrectChoice \texttt{@Bean}
  \choice \texttt{@Controller}
  \choice \texttt{@Service}
\end{choices}
\begin{solution}
  \textbf{Stereotypes} (\texttt{@Component}, \texttt{@Service}, \texttt{@Repository}, \texttt{@Controller}, \ldots) mark classes for component scanning. \texttt{@Bean} is used on \textbf{factory methods} (often inside \texttt{@Configuration}) to register a bean, not as a class-level stereotype.
\end{solution}

\question[4]
When requesting a bean with the scope \texttt{prototype}, you will always be given a new instance of the bean.
\begin{choices}
  \CorrectChoice true
  \choice false
\end{choices}
\begin{solution}
  \textbf{Prototype} scope returns a \textbf{new} instance for each \texttt{getBean} or injection point (depending on how it is resolved), unlike \textbf{singleton}, where one shared instance is typical.
\end{solution}

\question[4]
The default scope of a bean is \texttt{singleton}.
\begin{choices}
  \CorrectChoice true
  \choice false
\end{choices}
\begin{solution}
  In Spring, the default scope for beans in the IoC container is \textbf{singleton}: one logical instance per container (for that bean name).
\end{solution}

\question[4]
One of the problems Spring Core addresses is
\begin{choices}
  \CorrectChoice Tight coupling
  \choice Loose coupling
  \choice Pagination
  \choice Session management
\end{choices}
\begin{solution}
  Spring's \textbf{Inversion of Control} and \textbf{dependency injection} reduce \textbf{tight coupling} by injecting collaborators instead of hard-coding \texttt{new}. The goal is \emph{loose} coupling; pagination and sessions are separate concerns.
\end{solution}

\question[4]
The \texttt{@Configuration} annotation is used to
\begin{choices}
  \choice Provide XML configuration.
  \CorrectChoice Define a class as a configuration bean.
  \choice Provide annotation driven configuration.
  \choice Do the same as the \texttt{@Autowired} annotation.
\end{choices}
\begin{solution}
  \texttt{@Configuration} marks a class as a source of \textbf{@Bean} definitions (Java-based config). It does not replace all XML by itself; \texttt{@Autowired} declares dependency injection points, not configuration structure.
\end{solution}

\question[4]
You can provide configuration to Spring using XML, a configuration bean, and annotation-driven configuration.
\begin{choices}
  \CorrectChoice true
  \choice false
\end{choices}
\begin{solution}
  Spring supports \textbf{XML}, \textbf{Java config} (\texttt{@Configuration} classes), and \textbf{annotation-driven} setup (e.g.\ component scanning with \texttt{@Component}, \texttt{@Service}, etc.), often combined in one application.
\end{solution}

\question[4]
Which of the following dependency injection types are supported by the Spring Framework?
\begin{choices}
  \choice setter injection
  \choice constructor injection
  \CorrectChoice Both of the above.
  \choice None of the above.
\end{choices}
\begin{solution}
  Spring supports \textbf{constructor} and \textbf{setter} injection (and field injection via \texttt{@Autowired}). Interface injection is less central in modern Spring.
\end{solution}

\question[4]
Which of the following is an advantage of using Setter Injection in Spring?
\begin{choices}
  \choice It ensures that all dependencies are provided at the time of object creation.
  \CorrectChoice It allows for optional dependencies to be injected.
  \choice It allows for dependencies to be set using reflection.
  \choice None of the above
\end{choices}
\begin{solution}
  Setters allow \textbf{optional} collaborators to be omitted or changed after construction. \textbf{Constructor} injection is preferred for required, immutable dependencies.
\end{solution}

\question[4]
The \texttt{@RequestMapping} annotation is used to
\begin{choices}
  \CorrectChoice Provide a URL mapping to a resource for any HTTP method.
  \choice Provide a URL mapping to a resource for only GET methods.
  \choice Map a path variable to an array index.
  \choice Provide data contained in a request body as a \texttt{HashMap}.
\end{choices}
\begin{solution}
  \texttt{@RequestMapping} can map a path (and optionally HTTP methods via \texttt{method = \ldots}). Narrower annotations (\texttt{@GetMapping}, \texttt{@PostMapping}, \ldots) are specializations. Path variables use \texttt{@PathVariable}; the body is bound with \texttt{@RequestBody} or parameters, not \texttt{@RequestMapping} alone.
\end{solution}

\question[4]
What is the \texttt{@Controller} annotation?
\begin{choices}
  \CorrectChoice The \texttt{@Controller} annotation indicates that a particular class serves the role of a controller.
  \choice The \texttt{@Controller} annotation indicates how to control the transaction management.
  \choice The \texttt{@Controller} annotation indicates how to control the dependency injection.
  \choice The \texttt{@Controller} annotation indicates how to control the aspect programming.
\end{choices}
\begin{solution}
  In Spring MVC, \texttt{@Controller} (stereotype of \texttt{@Component}) marks a web controller whose methods handle HTTP via \texttt{@RequestMapping} and related annotations. It does not govern transactions, DI, or AOP by itself.
\end{solution}

\question[4]
When used on the class, the \texttt{@RestController} annotation implies what on all methods?
\begin{choices}
  \CorrectChoice \texttt{@ResponseBody}
  \choice \texttt{@RequestBody}
  \choice \texttt{@RequestMapping}
  \choice \texttt{@RequestParam}
\end{choices}
\begin{solution}
  \texttt{@RestController} is \texttt{@Controller} + \texttt{@ResponseBody}: return values of handler methods are serialized to the HTTP response body (e.g.\ JSON) by default. \texttt{@RequestBody} binds incoming body; \texttt{@RequestMapping} maps URLs; \texttt{@RequestParam} binds query/form params.
\end{solution}

\question[4]
The \texttt{@RequestBody} annotation is used to
\begin{choices}
  \choice Define a class as a controller.
  \choice Define that the return value of a method should be serialized to some form such as XML, JSON, ect.
  \CorrectChoice Define that the request body should be deserialized to a Java object from some form such as XML, JSON, ect.
  \choice Parse a variable in the requested URL and inject it into the method call as a parameter.
\end{choices}
\begin{solution}
  \texttt{@RequestBody} binds the \textbf{HTTP request body} to a Java object using message converters (often JSON). Response serialization uses \texttt{@ResponseBody} or \texttt{@RestController}. (Original uses \textit{ect.})
\end{solution}

\question[4]
The \texttt{@PathVariable} annotation is used to
\begin{choices}
  \choice Define a class as a controller.
  \choice Define that the return value of a method should be serialized to some form such as XML, JSON, ect.
  \choice Define that the request body should be deserialized to a Java object from form such as XML, JSON, ect.
  \CorrectChoice Parse the variable in the requested URL and provide it to the method call as a parameter.
\end{choices}
\begin{solution}
  \texttt{@PathVariable} binds URI template variables (e.g.\ \texttt{/orders/\{id\}}) to method parameters. It is not for the request body or return serialization.
\end{solution}

\question[4]
Spring MVC. Given the following annotation, what would be the most appropriate handler syntax?
\texttt{@GetMapping("/\{id\}")}
\begin{choices}
  \choice \texttt{public Object getById(@PathParam(name="id") int value) \{\}}
  \choice \texttt{public Object getById(@RequestBody int id) \{\}};
  \choice \texttt{public Object getById(@RequestMapping int id) \{\}}
  \CorrectChoice \texttt{public Object getById(@PathVariable int id) \{\}}
\end{choices}
\begin{solution}
  Path segments like \texttt{\{id\}} bind with Spring's \texttt{@PathVariable}. \texttt{@PathParam} is JAX-RS, not Spring MVC. \texttt{@RequestBody} is for the HTTP body; \texttt{@RequestMapping} is not a parameter annotation.
\end{solution}

\question[4]
The \texttt{@Transactional} annotation is used to
\begin{choices}
  \CorrectChoice Define that the underlying code should be executed as a database transaction.
  \choice Define that the underlying code is used to make a purchase.
  \choice Define that the underlying code should never throw an exception.
  \choice Does nothing.
\end{choices}
\begin{solution}
  \texttt{@Transactional} (Spring's declarative transaction support) demarcates a \textbf{transaction} boundary, typically around service methods hitting a transactional resource. Behavior on exceptions depends on rollback rules.
\end{solution}

\question[4]
Spring Data's \texttt{JpaRepository} features a fluent API which allows it to
\begin{choices}
  \choice Generate Procedural Langauge functions on the database.
  \choice Know whether data exists without checking the database.
  \choice Provide an in-memory database that doesn't require a separate database.
  \CorrectChoice Generate SQL queries by parsing an abstract method's name.
\end{choices}
\begin{solution}
  \textbf{Query methods} (derived from method names like \texttt{findByLastName}) let Spring Data JPA build queries from signatures. The first choice keeps the original typo \textit{Langauge}. Existence still hits the database (e.g.\ \texttt{existsBy\ldots}); in-memory DB is not what \texttt{JpaRepository} itself provides.
\end{solution}

\question[4]
Regarding Spring Data, what are custom query methods?
\begin{choices}
  \choice Methods whose names define what query to perform using property expressions. You will need to create a class that implements these methods.
  \CorrectChoice Methods whose names define what query to perform using property expressions. Spring will automatically generate the implementation for you.
  \choice Methods that allow developers to write SQL queries directly in Java code.
  \choice Predefined methods provided by Spring Data for common database operations.
\end{choices}
\begin{solution}
  \textbf{Derived query methods} on a repository interface are implemented by the Spring Data infrastructure at runtime from the method name and domain model. You can also use \texttt{@Query} for explicit JPQL/SQL.
\end{solution}

\uplevel{\par\textbf{Spring Data annotations} (from the fill-in table in the source bank).\par}

\question[4]
Which annotation marks a Spring Data repository \emph{interface} that should \textbf{not} be picked up as a Spring Data repository implementation (for example a shared base interface with only common signatures)?
\begin{choices}
  \CorrectChoice \texttt{@NoRepositoryBean}
  \choice \texttt{@Repository}
  \choice \texttt{@Query}
  \choice \texttt{@Id}
\end{choices}
\begin{solution}
  \texttt{@NoRepositoryBean} tells Spring Data not to create a proxy for that interface---useful for intermediate templates. Concrete repository interfaces (extending \texttt{JpaRepository}, etc.) are normally discovered without this marker.
\end{solution}

\question[4]
Which annotation marks a field in a JPA entity class as the primary key? (The source table had a typo \texttt{@ld}.)
\begin{choices}
  \choice \texttt{@NoRepositoryBean}
  \choice \texttt{@ld}
  \CorrectChoice \texttt{@Id}
  \choice \texttt{@Transient}
\end{choices}
\begin{solution}
  JPA uses \texttt{@Id} on the primary-key field or property. \texttt{@ld} is a typo in the original key column; \texttt{@Transient} means not persisted.
\end{solution}

\question[4]
Which annotation marks a field to be ignored by the persistence provider during reads and writes (not stored in the database)?
\begin{choices}
  \choice \texttt{@Id}
  \choice \texttt{@Query}
  \CorrectChoice \texttt{@Transient}
  \choice \texttt{@NoRepositoryBean}
\end{choices}
\begin{solution}
  \texttt{javax.persistence.Transient} (or \texttt{jakarta.persistence.Transient}) excludes the field from the entity mapping. Spring Data MongoDB also has field-level transient support in its mapping model.
\end{solution}

\question[4]
Which annotation supplies a JPQL (or native) query string for a Spring Data repository method?
\begin{choices}
  \choice \texttt{@NoRepositoryBean}
  \choice \texttt{@Id}
  \choice \texttt{@Transient}
  \CorrectChoice \texttt{@Query}
\end{choices}
\begin{solution}
  \texttt{@Query} on a repository method declares JPQL/SQL explicitly when derived names are insufficient. Parameters can be bound by position or \texttt{@Param}.
\end{solution}

\question[4]
To use the Spring framework you must use Spring Boot
\begin{choices}
  \choice TRUE
  \CorrectChoice FALSE
\end{choices}
\begin{solution}
  The \textbf{Spring Framework} predates \textbf{Spring Boot}. Boot is a convention-over-configuration layer on top; plain Spring with XML or Java config remains valid without Boot.
\end{solution}

\question[4]
What is the purpose of using \texttt{MockMvc} for Spring MVC testing?
\begin{choices}
  \choice It is a replacement for JUnit and Mockito
  \CorrectChoice It can run integration tests that test controllers without explicitly starting a Servlet container
  \choice It can run regression tests that detect defects introduced by new code
  \choice It enables making calls to external application's REST API endpoints for testing
\end{choices}
\begin{solution}
  \texttt{MockMvc} (from \texttt{spring-test}) dispatches through the \textbf{DispatcherServlet} test infrastructure, so you exercise controller mapping, validation, and serialization without binding to a live port. It is often used from Spring Boot tests but is not a Boot-only API. JUnit/Mockito remain separate tools.
\end{solution}

\uplevel{\par\textbf{Spring Framework practice bank}\par}

\question[4]
Which statement best describes constructor injection in Spring?
\begin{choices}
  \CorrectChoice It makes required dependencies explicit at object creation and is generally preferred for mandatory collaborators
  \choice It is only available when XML configuration is used
  \choice It can only inject primitive values, not beans
  \choice It is deprecated in favor of field injection
\end{choices}
\begin{solution}
  Constructor injection works well for required dependencies and supports immutability/testing. Field injection still exists but is less explicit and harder to test in isolation.
\end{solution}

\question[4]
What is the most accurate difference between \texttt{@Controller} and \texttt{@RestController}?
\begin{choices}
  \choice \texttt{@Controller} cannot map HTTP requests
  \CorrectChoice \texttt{@RestController} implies \texttt{@ResponseBody} semantics on handler methods, while \texttt{@Controller} is typically used for MVC view rendering unless methods are annotated otherwise
  \choice \texttt{@RestController} is only valid in Spring Boot, not Spring Framework
  \choice They are aliases with no behavioral differences
\end{choices}
\begin{solution}
  \texttt{@RestController} combines \texttt{@Controller} and \texttt{@ResponseBody}, so return values are serialized to the response body by default.
\end{solution}

\question[4]
Given \texttt{@GetMapping("/orders/\{id\}")}, which parameter binding is correct?
\begin{choices}
  \choice \texttt{@RequestBody Long id}
  \CorrectChoice \texttt{@PathVariable Long id}
  \choice \texttt{@RequestParam Long id}
  \choice \texttt{@Autowired Long id}
\end{choices}
\begin{solution}
  URI template tokens (\texttt{\{id\}}) are bound with \texttt{@PathVariable}; \texttt{@RequestParam} binds query string/form parameters.
\end{solution}

\question[4]
Which statement about \texttt{@Transactional} default rollback behavior is correct?
\begin{choices}
  \CorrectChoice By default, Spring rolls back on unchecked exceptions (\texttt{RuntimeException} and \texttt{Error}) and commits on checked exceptions unless configured otherwise
  \choice It always rolls back for every thrown exception type
  \choice It only applies when methods are static
  \choice It opens one global JVM-wide transaction shared by all beans
\end{choices}
\begin{solution}
  You can customize rollback rules (e.g.\ \texttt{rollbackFor = Exception.class}) when checked exception rollback is needed.
\end{solution}

\question[4]
Which Spring Data repository method name will derive a query for users by exact email?
\begin{choices}
  \choice \texttt{queryUserEmail(String email)}
  \CorrectChoice \texttt{findByEmail(String email)}
  \choice \texttt{selectEmailFromUser(String email)}
  \choice \texttt{getEmailWhereUser(String email)}
\end{choices}
\begin{solution}
  Spring Data parses supported keywords/patterns (like \texttt{findBy...}, \texttt{existsBy...}) to generate query implementations.
\end{solution}

\question[4]
Which annotation should you use to inject a query parameter like \texttt{?page=2} into a controller method argument?
\begin{choices}
  \CorrectChoice \texttt{@RequestParam}
  \choice \texttt{@PathVariable}
  \choice \texttt{@RequestBody}
  \choice \texttt{@Configuration}
\end{choices}
\begin{solution}
  Query string values are bound via \texttt{@RequestParam}; body payloads use \texttt{@RequestBody}; path segments use \texttt{@PathVariable}.
\end{solution}

\uplevel{\par\textbf{Additional Spring Framework Concepts: Exception Handling}\par}

\question[4]
Which annotation is used to define a centralized, global exception handler for REST APIs, automatically serializing the return value to the HTTP response body?
\begin{choices}
  \choice \texttt{@ControllerAdvice}
  \CorrectChoice \texttt{@RestControllerAdvice}
  \choice \texttt{@ExceptionHandler}
  \choice \texttt{@ResponseStatus}
\end{choices}
\begin{solution}
  \texttt{@RestControllerAdvice} combines \texttt{@ControllerAdvice} and \texttt{@ResponseBody}. It is used to handle exceptions globally across all \texttt{@RestController} components and write the result directly to the response body.
\end{solution}

\question[4]
Which annotation is used to define a centralized, global exception handler typically for classic Spring MVC applications where the handler methods might return view names?
\begin{choices}
  \CorrectChoice \texttt{@ControllerAdvice}
  \choice \texttt{@RestControllerAdvice}
  \choice \texttt{@ExceptionHandler}
  \choice \texttt{@ResponseStatus}
\end{choices}
\begin{solution}
  \texttt{@ControllerAdvice} intercepts exceptions globally across \texttt{@Controller} classes. Unlike \texttt{@RestControllerAdvice}, it does not imply \texttt{@ResponseBody} by default.
\end{solution}

\question[4]
Inside an advice class or a controller, which annotation is placed directly on a method to indicate that the method should be invoked when a specific exception type is thrown?
\begin{choices}
  \choice \texttt{@ControllerAdvice}
  \choice \texttt{@RestControllerAdvice}
  \CorrectChoice \texttt{@ExceptionHandler}
  \choice \texttt{@ResponseStatus}
\end{choices}
\begin{solution}
  The \texttt{@ExceptionHandler} annotation marks a specific method to handle a specific exception (e.g., \texttt{@ExceptionHandler(ClientValidationException.class)}).
\end{solution}

\question[4]
Which annotation can be applied to an exception class or an exception handler method to explicitly set the HTTP status code (e.g., 404 Not Found) that should be returned?
\begin{choices}
  \choice \texttt{@ControllerAdvice}
  \choice \texttt{@RestControllerAdvice}
  \choice \texttt{@ExceptionHandler}
  \CorrectChoice \texttt{@ResponseStatus}
\end{choices}
\begin{solution}
  \texttt{@ResponseStatus} dictates the exact HTTP status code and an optional reason phrase to be returned when the exception is thrown or the handler method is executed.
\end{solution}

\uplevel{\par\textbf{Additional Spring Framework Concepts: REST Mappings}\par}

\question[4]
Which Spring Web annotation maps HTTP GET requests to a specific handler method, typically used to retrieve a resource without modifying server state?
\begin{choices}
  \CorrectChoice \texttt{@GetMapping}
  \choice \texttt{@PostMapping}
  \choice \texttt{@DeleteMapping}
  \choice \texttt{@PatchMapping}
\end{choices}
\begin{solution}
  \texttt{@GetMapping} is a shortcut for \texttt{@RequestMapping(method = RequestMethod.GET)}. It is used for read-only operations.
\end{solution}

\question[4]
Which Spring Web annotation maps HTTP POST requests to a specific handler method, typically used to create a new resource or submit data to the server?
\begin{choices}
  \choice \texttt{@GetMapping}
  \CorrectChoice \texttt{@PostMapping}
  \choice \texttt{@DeleteMapping}
  \choice \texttt{@PatchMapping}
\end{choices}
\begin{solution}
  \texttt{@PostMapping} is a shortcut for \texttt{@RequestMapping(method = RequestMethod.POST)}. It is used when an operation is not idempotent, like creating a new user or message.
\end{solution}

\question[4]
Which Spring Web annotation maps HTTP DELETE requests to a specific handler method, typically used to remove an existing resource?
\begin{choices}
  \choice \texttt{@GetMapping}
  \choice \texttt{@PostMapping}
  \CorrectChoice \texttt{@DeleteMapping}
  \choice \texttt{@PatchMapping}
\end{choices}
\begin{solution}
  \texttt{@DeleteMapping} is a shortcut for \texttt{@RequestMapping(method = RequestMethod.DELETE)}. It maps URL paths to a deletion operation.
\end{solution}

\question[4]
Which Spring Web annotation maps HTTP PATCH requests to a specific handler method, typically used to apply partial modifications to an existing resource?
\begin{choices}
  \choice \texttt{@GetMapping}
  \choice \texttt{@PostMapping}
  \choice \texttt{@DeleteMapping}
  \CorrectChoice \texttt{@PatchMapping}
\end{choices}
\begin{solution}
  \texttt{@PatchMapping} is a shortcut for \texttt{@RequestMapping(method = RequestMethod.PATCH)}. It is idiomatically used for partial updates, like changing just the text of a message.
\end{solution}

\uplevel{\par\textbf{Additional Spring Framework Concepts: Data and Architecture}\par}

\question[4]
In a layered architecture, what is the primary purpose of the \texttt{@Service} annotation?
\begin{choices}
  \choice It exposes HTTP endpoints to the web.
  \choice It provides direct access to the database using SQL.
  \CorrectChoice It marks a class that holds business logic and coordinates operations between the controller and repository layers.
  \choice It automatically manages user sessions and cookies.
\end{choices}
\begin{solution}
  The \texttt{@Service} stereotype annotation is used to designate a class as a service layer component, encapsulating business logic and validation.
\end{solution}

\question[4]
What is the primary benefit of returning a \texttt{ResponseEntity} from a Spring MVC controller method instead of directly returning the domain object?
\begin{choices}
  \choice It automatically serializes the object to XML.
  \choice It allows the controller to handle transactions automatically.
  \CorrectChoice It allows you to fully configure the entire HTTP response, including the status code, headers, and body.
  \choice It enforces security and prevents unauthorized access.
\end{choices}
\begin{solution}
  \texttt{ResponseEntity} represents the entire HTTP response. Returning it gives developers fine-grained control over the HTTP status code (e.g., 200 OK, 400 Bad Request, 409 Conflict), headers, and the response body.
\end{solution}

\question[4]
When defining a repository interface by extending \texttt{JpaRepository<T, ID>}, what do the generic type parameters \texttt{T} and \texttt{ID} represent?
\begin{choices}
  \choice \texttt{T} is the table name, and \texttt{ID} is the index.
  \CorrectChoice \texttt{T} is the domain entity class, and \texttt{ID} is the data type of the entity's primary key.
  \choice \texttt{T} is the transaction manager, and \texttt{ID} is the database connection ID.
  \choice \texttt{T} is the return type of the queries, and \texttt{ID} is the query string.
\end{choices}
\begin{solution}
  In \texttt{JpaRepository<T, ID>}, \texttt{T} specifies the JPA entity type the repository manages, and \texttt{ID} specifies the Java type of that entity's \texttt{@Id} field.
\end{solution}

\question[4]
In a JPA entity class, what is the purpose of the \texttt{@GeneratedValue} annotation when applied alongside \texttt{@Id}?
\begin{choices}
  \choice It generates a random UUID for password hashing.
  \CorrectChoice It configures the way the primary key value is automatically generated by the database.
  \choice It ensures that the field is never serialized into JSON.
  \choice It automatically creates a new table in the database if one does not exist.
\end{choices}
\begin{solution}
  \texttt{@GeneratedValue} specifies the generation strategy for primary keys (like auto-increment columns), allowing the database or persistence provider to automatically assign an ID when a new entity is saved.
\end{solution}

\uplevel{\par\textbf{Additional practice: Spring Framework (polished Q\&A bank)}\par}

\question[4]
What is the purpose of \texttt{ResponseEntity} in the following \texttt{RestTemplate} example?
\begin{lstlisting}[style=java]
RestTemplate rest = new RestTemplate();
HttpEntity<DataModel> request = new HttpEntity<>(new DataModel("value"));
String resourceUrl = "http://localhost:8080/myservice/data";
ResponseEntity<DataModel> response =
    rest.exchange(resourceUrl, HttpMethod.POST, request, DataModel.class);
\end{lstlisting}
\begin{choices}
  \choice \texttt{ResponseEntity} represents only the HTTP response body.
  \choice \texttt{ResponseEntity} specifies only the expected Java type of the request entity.
  \choice It is only used to set custom HTTP headers on the outgoing request.
  \CorrectChoice It represents the full HTTP response (status, headers, and body), and the type parameter describes the deserialized response body.
\end{choices}
\begin{solution}
  \texttt{ResponseEntity<T>} wraps the complete client-side (or server-side) HTTP outcome: status line, headers, and typed body. The \texttt{Class}/\texttt{ParameterizedTypeReference} argument to \texttt{exchange} tells \texttt{RestTemplate} how to convert the response body.
\end{solution}

\question[4]
When injecting property placeholders from a file such as \texttt{jdbc.properties} into fields in a \texttt{@Configuration} class, which annotation fills the blank?
\begin{lstlisting}[style=java]
@Configuration
@ImportResource("classpath:/com/example/resources/properties-config.xml")
public class AppConfig {

    @________("${jdbc.url}")
    private String url;

    @________("${jdbc.username}")
    private String username;

    @________("${jdbc.password}")
    private String password;

    @Bean
    public DataSource dataSource() {
        return new DriverManagerDataSource(url, username, password);
    }
}
\end{lstlisting}
\begin{choices}
  \CorrectChoice \texttt{@Value}
  \choice \texttt{@Inject}
  \choice \texttt{@Qualifier}
  \choice \texttt{@Resource}
\end{choices}
\begin{solution}
  \texttt{@Value("\${...}")} injects scalar values from the environment, property sources, or SpEL. \texttt{@Qualifier} disambiguates beans by name; \texttt{@Resource} is the JSR-250 annotation; \texttt{@Inject} is the JSR-330 counterpart to \texttt{@Autowired}.
\end{solution}

\question[4]
Which XML snippet correctly performs \textbf{setter injection} of \texttt{orderDao} into \texttt{OrderServiceImpl}?
\begin{choices}
  \CorrectChoice \texttt{<bean id="orderServ" class="com.example.services.OrderServiceImpl"><property name="orderDao" ref="orderDao"/></bean>}
  \choice \texttt{<bean id="orderServ" class="com.example.services.OrderServiceImpl"></bean>}
  \choice Both snippets are equally correct for setter injection.
  \choice Neither snippet can work in Spring XML wiring.
\end{choices}
\begin{solution}
  The \texttt{<property name="..." ref="..."/>} form calls the matching \texttt{setOrderDao} mutator (by JavaBeans conventions). An empty \texttt{<bean>} can still construct the object but does not inject collaborators.
\end{solution}

\question[4]
Which injection style is most appropriate when you want minimal boilerplate and are least concerned about testability and explicit dependencies?
\begin{choices}
  \CorrectChoice Field injection
  \choice Setter injection
  \choice Constructor injection
  \choice All of the above are equivalent for that goal
\end{choices}
\begin{solution}
  Field injection (\texttt{@Autowired} on fields) is terse but hides dependencies, complicates unit tests, and makes immutability harder. Constructor injection remains the default recommendation for required dependencies.
\end{solution}

\question[4]
Explain the transaction semantics of the method below.
\begin{lstlisting}[style=java]
@Service
@Transactional
public class MyService {
    @Transactional(
        isolation = Isolation.SERIALIZABLE,
        rollbackFor = { ConstraintViolationException.class },
        propagation = Propagation.REQUIRES_NEW)
    public Integer createNewData(...) { ... }
}
\end{lstlisting}
\begin{choices}
  \CorrectChoice A new transaction starts for \texttt{createNewData}; an existing transaction is suspended; \texttt{ConstraintViolationException} triggers rollback.
  \choice A new transaction starts only if none exists; otherwise the call throws immediately.
  \choice The annotations create a JVM checkpoint so failures always resume the prior program counter.
  \choice Rollback applies to every checked exception except \texttt{ConstraintViolationException}.
\end{choices}
\begin{solution}
  \texttt{Propagation.REQUIRES\_NEW} always opens a new physical transaction and suspends the surrounding one if present. \texttt{rollbackFor} extends rollback behavior to the listed checked exceptions. \texttt{SERIALIZABLE} is the strictest isolation level.
\end{solution}

\question[4]
Consider a controller declared as follows. Which statement is true, and what single-annotation alternative preserves the behavior?
\begin{lstlisting}[style=java]
@Controller
@ResponseBody
public class MyController { ... }
\end{lstlisting}
\begin{choices}
  \CorrectChoice Handler return values are written to the HTTP response body; \texttt{@RestController} can replace both class-level annotations.
  \choice Handler return values are written to the HTTP response body; there is no equivalent composed annotation.
  \choice Return values always select a view name; replace \texttt{@ResponseBody} with \texttt{@View}.
  \choice The combination is invalid in Spring MVC.
\end{choices}
\begin{solution}
  Class-level \texttt{@ResponseBody} applies to every handler method on the controller. \texttt{@RestController} is meta-annotated with \texttt{@Controller} and \texttt{@ResponseBody}, so it is the idiomatic replacement.
\end{solution}

\question[4]
What does this code accomplish?
\begin{lstlisting}[style=java]
@RestController
@RequestMapping("/api")
public class API {
    @GetMapping
    public String getData() {
        return "Hello from API!";
    }
}
\end{lstlisting}
\begin{choices}
  \choice It issues an outbound GET request automatically at startup.
  \CorrectChoice It exposes an HTTP GET endpoint (by default on \texttt{/api}) that returns the given \texttt{String} body.
  \choice It registers a POST handler because \texttt{@GetMapping} without a path implies POST.
  \choice \texttt{@RequestMapping} cannot be applied at class level.
\end{choices}
\begin{solution}
  Class-level \texttt{@RequestMapping("/api")} prefixes member mappings; a method-level \texttt{@GetMapping} without a path maps to the class path, so clients \texttt{GET /api} receive the literal greeting (subject to content negotiation defaults).
\end{solution}

\question[4]
In the controller fragment below, what is the URI template variable name, and what does the code demonstrate?
\begin{lstlisting}[style=java]
@Controller
@RequestMapping("my_controller")
@ResponseBody
public class MyController {
    @GetMapping(path = "my_data/{id}")
    public ResponseEntity<?> getDataById(
            @PathVariable(name = "id", required = false) Integer dataId) {
        if (dataId != null) {
            // id segment present
        } else {
            // id omitted or not parseable
        }
        return ResponseEntity.ok().build();
    }
}
\end{lstlisting}
\begin{choices}
  \CorrectChoice The template variable is \texttt{id}; the class is mapped under \texttt{/my\_controller} and the handler under \texttt{/my\_data/\{id\}}; \texttt{@PathVariable} binds \texttt{id} to \texttt{dataId}.
  \choice The template variable is \texttt{dataId}; \texttt{@GetMapping} supplies the base path for every CRUD verb automatically.
  \choice The template variable is \texttt{id}; \texttt{@GetMapping} alone defines the entire servlet path without class-level mapping.
  \choice The template variable is \texttt{dataId}; \texttt{@PathVariable} matches only the \texttt{type} attribute, not \texttt{name}.
\end{choices}
\begin{solution}
  The brace name in \texttt{\{id\}} is the path variable. The method parameter may use any Java identifier; \texttt{name = "id"} binds it to that segment. \texttt{required = false} allows missing segments in some routing setups, though a missing template segment usually yields no match---the keyed teaching item stresses naming vs.\ parameter identifier.
\end{solution}

\question[4]
Given this bean declaration, which line correctly retrieves the bean from an \texttt{ApplicationContext} named \texttt{context}?
\begin{lstlisting}[style=java]
@Bean(name = "myBeanName")
public Student student() {
    return new Student();
}
\end{lstlisting}
\begin{choices}
  \CorrectChoice \texttt{Student s = context.getBean("myBeanName", Student.class);}
  \choice \texttt{Student s = (Student) context.getBean("student");}
  \choice \texttt{Student s = context.getBean("Student", Student.class);}
  \choice \texttt{Student s = (Student) context.getBean("mybeanname");}
\end{choices}
\begin{solution}
  Explicit bean names are case-sensitive; \texttt{"myBeanName"} matches the \texttt{@Bean} name. Prefer the typed overload \texttt{getBean(String, Class)} over casting. The default method name \texttt{student} would register \texttt{"student"} only if no \texttt{name} were given.
\end{solution}

\question[4]
\textit{(Select all that apply.)} What are some ways to resolve ambiguity when Spring wires constructor arguments for a \texttt{@Bean} method or component?
\begin{checkboxes}
  \CorrectChoice Index (argument position when combined with other disambiguators)
  \CorrectChoice Type (distinct parameter types)
  \CorrectChoice Parameter name (retained names or explicit \texttt{@Qualifier} metadata)
  \choice \texttt{Value} alone without qualifiers
\end{checkboxes}
\begin{solution}
  Spring uses types, ordering metadata, qualifiers, and parameter names (when available) to select collaborators. A bare \texttt{@Value} does not address multiple bean candidates of the same type.
\end{solution}

\question[4]
What does the following Spring Data JPA arrangement accomplish?
\begin{lstlisting}[style=java]
public interface MyRepoBase<T, ID> extends JpaRepository<T, ID> {
    List<T> findAllByPrice(Double price);
}

@Repository
public interface BookRepo extends MyRepoBase<Book, Long> { }
\end{lstlisting}
\begin{choices}
  \CorrectChoice \texttt{MyRepoBase} is a reusable contract; Spring creates a proxy only for concrete repository interfaces such as \texttt{BookRepo}, which inherits \texttt{JpaRepository} and \texttt{findAllByPrice}.
  \choice Spring instantiates both \texttt{MyRepoBase} and \texttt{BookRepo} as separate beans.
  \choice Spring ignores \texttt{BookRepo} and only creates \texttt{MyRepoBase}.
  \choice None of the above.
\end{choices}
\begin{solution}
  Intermediate interfaces are not beans unless annotated with \texttt{@NoRepositoryBean}. Declared repositories (\texttt{BookRepo}) receive generated implementations that include custom finder methods declared on shared parents.
\end{solution}

\question[4]
Why is \texttt{@Transactional} commonly applied at the \texttt{@Service} layer instead of on entities or controllers alone?
\begin{choices}
  \CorrectChoice A transaction encapsulates application-level use cases that may span multiple repositories and rules, not just persistence mapping.
  \choice Transactions should be tied exclusively to JPA entity lifecycle callbacks.
  \choice Services are the only layer allowed to open JDBC connections.
  \choice Controllers cannot participate in Spring's transaction proxy infrastructure.
\end{choices}
\begin{solution}
  Services orchestrate domain operations and coordinate repositories; that is the natural boundary for atomic business transactions. Controllers should stay thin; entities map state, not workflow.
\end{solution}

\question[4]
Why is specifying \texttt{defaultValue} on \texttt{@RequestParam} useful?
\begin{lstlisting}[style=java]
@GetMapping(path = "my_data")
public ResponseEntity<?> getData(
        @RequestParam(name = "tag", required = false, defaultValue = "none") String tag) {
    ...
}
\end{lstlisting}
\begin{choices}
  \CorrectChoice When the query parameter is omitted, Spring still binds the argument using the supplied default string after type conversion.
  \choice It lazily constructs domain objects without any HTTP interaction.
  \choice It forces every client to send identical query strings regardless of input.
  \choice It replaces the need for \texttt{required = false}.
\end{choices}
\begin{solution}
  \texttt{defaultValue} implies \texttt{required = false} in typical usage: missing parameters fall back to the literal string before conversion to the target type. Values must be compile-time constant strings; dynamic defaults belong in method logic.
\end{solution}

\question[4]
\textit{(Select all that apply.)} Which items are part of the Spring \textbf{Core} container story (as commonly grouped in courses)?
\begin{checkboxes}
  \CorrectChoice Spring IoC container
  \choice Model--View--Controller support (\texttt{DispatcherServlet})
  \CorrectChoice \texttt{ApplicationContext} as the richer \texttt{BeanFactory}
  \choice JDBC template module by itself
  \CorrectChoice Managed Spring beans
\end{checkboxes}
\begin{solution}
  Core focuses on the IoC container, bean lifecycle, and \texttt{ApplicationContext}. MVC lives in Spring Web; JDBC support is its own module even though apps combine them constantly.
\end{solution}

\question[4]
Which option is \textbf{not} a benefit of adopting the Spring Framework?
\begin{choices}
  \choice Modularity across starters and modules
  \choice Emphasis on plain Java objects for business code
  \choice Inversion of Control and declarative middleware
  \CorrectChoice ``Complexity of the framework'' as an inherent benefit
\end{choices}
\begin{solution}
  Spring's breadth can add learning curve, but that complexity is a trade-off, not a marketed advantage. Modularity, POJO-centric design, and IoC are the usual positives.
\end{solution}

\question[4]
\textit{(Select all that apply.)} Which styles are common Spring dependency injection mechanisms?
\begin{checkboxes}
  \CorrectChoice Constructor injection
  \choice Object injection (not a Spring term)
  \choice Container injection (not a Spring term)
  \CorrectChoice Setter injection
\end{checkboxes}
\begin{solution}
  Spring supports constructor and setter injection as first-class patterns, plus field injection via \texttt{@Autowired}. ``Object injection'' and ``container injection'' are distractors.
\end{solution}

\question[4]
What kind of injection does this XML fragment configure?
\begin{lstlisting}[style=java]
<property name="orderDao" ref="orderDaoBean"/>
\end{lstlisting}
\begin{choices}
  \CorrectChoice Setter injection
  \choice Property injection (distinct from Spring's setter wiring)
  \choice Constructor injection
  \choice Container injection
\end{choices}
\begin{solution}
  The \texttt{<property>} element invokes JavaBeans-style setters on the bean instance after construction.
\end{solution}

\question[4]
What kind of injection does this XML fragment configure?
\begin{lstlisting}[style=java]
<constructor-arg name="orderDao" ref="orderDaoBean"/>
\end{lstlisting}
\begin{choices}
  \choice Property injection
  \CorrectChoice Constructor injection
  \choice Setter injection
  \choice Object injection
\end{choices}
\begin{solution}
  \texttt{<constructor-arg>} supplies arguments to a chosen constructor or factory method, which is constructor-based DI.
\end{solution}

\question[4]
What does the \texttt{@Repository} stereotype communicate?
\begin{choices}
  \CorrectChoice It is a \texttt{@Component} specialization marking persistence/dao-style components and enables exception translation.
  \choice It is primarily for Spring MVC handler mapping alongside \texttt{@RequestMapping}.
  \choice It marks service-layer orchestration beans.
  \choice It replaces \texttt{@Component} for every bean in the application.
\end{choices}
\begin{solution}
  \texttt{@Repository} signals data access responsibilities, enables classpath scanning, and (with persistence providers) supports Spring's \texttt{DataAccessException} translation machinery.
\end{solution}

\question[4]
What does the \texttt{@Controller} stereotype communicate?
\begin{choices}
  \choice It is a DAO specialization for CRUD operations.
  \CorrectChoice It marks a Spring MVC controller whose \texttt{@RequestMapping}-style methods are dispatched by \texttt{DispatcherServlet}.
  \choice It marks transactional service facades.
  \choice It is unrelated to web request handling.
\end{choices}
\begin{solution}
  \texttt{@Controller} is a web stereotype discovered by component scanning; handler methods map HTTP to service calls or view names.
\end{solution}

\question[4]
Which annotations are Spring stereotype markers at the component level?
\begin{choices}
  \choice Only \texttt{@Service}
  \choice Only \texttt{@Controller} and \texttt{@Repository}
  \choice Only \texttt{@Component}
  \CorrectChoice \texttt{@Service}, \texttt{@Controller}, \texttt{@Repository}, and \texttt{@Component}
\end{choices}
\begin{solution}
  Each stereotype is meta-annotated with \texttt{@Component} (directly or transitively), enabling shared scanning semantics with specialized semantics per alias.
\end{solution}

\question[4]
In MVC, what role does the \textbf{Model} play?
\begin{choices}
  \choice It is the front controller that dispatches requests.
  \CorrectChoice It captures domain data and business rules the user cares about.
  \choice It owns view rendering and templates exclusively.
  \choice It terminates TCP connections for HTTP keep-alive.
\end{choices}
\begin{solution}
  The model is state and behavior the application manipulates; controllers update it and views (or serializers) present it.
\end{solution}

\question[4]
In MVC, what role does the \textbf{View} play?
\begin{choices}
  \choice It maps URLs to controller methods.
  \choice It stores transactional isolation levels.
  \CorrectChoice It presents model data to users (HTML templates, PDFs, JSON serializers in API-style MVC, etc.).
  \choice It replaces the need for controllers entirely.
\end{choices}
\begin{solution}
  Views render output for clients. In REST-style apps, the ``view'' is often the serialized representation produced by message converters.
\end{solution}

\question[4]
In MVC, what role does the \textbf{Controller} play?
\begin{choices}
  \CorrectChoice It translates incoming requests into operations on the model and selects the appropriate response or view.
  \choice It defines database indexes and migration scripts.
  \choice It is only responsible for static resource caching headers.
  \choice It replaces the model with DTOs in every application.
\end{choices}
\begin{solution}
  Controllers orchestrate request flow, delegate to services, and return views or bodies without embedding all business rules.
\end{solution}

\question[4]
In Spring MVC, what is the \textbf{DispatcherServlet}?
\begin{choices}
  \choice A JDBC connection pool implementation.
  \choice A JPA \texttt{EntityManagerFactory} facade.
  \CorrectChoice The front controller servlet that routes HTTP requests to Spring-managed handlers.
  \choice A Tomcat valve for gzip compression only.
\end{choices}
\begin{solution}
  \texttt{DispatcherServlet} is the central servlet entry point that consults handler mappings, invokes controllers, and applies view resolution or message conversion.
\end{solution}

\question[4]
Regarding ACID transactions, what does \textbf{Atomicity} guarantee?
\begin{choices}
  \choice The database moves between consistent states obeying all constraints at every intermediate statement.
  \choice Committed work survives power loss.
  \CorrectChoice Either all operations in the transaction complete successfully or none of their effects remain.
  \choice Concurrent transactions cannot observe each other's partial results.
\end{choices}
\begin{solution}
  Atomicity is the ``all or nothing'' guarantee for the transaction's operations.
\end{solution}

\question[4]
Regarding ACID transactions, what does \textbf{Consistency} mean?
\begin{choices}
  \CorrectChoice A successful transaction moves the database from one valid state to another valid state with respect to declared rules.
  \choice Committed work survives power loss.
  \choice Either every sub-operation completes or none do.
  \choice Transactions are hidden from one another.
\end{choices}
\begin{solution}
  Consistency ties to integrity constraints and business rules---illegal states cannot be committed.
\end{solution}

\question[4]
Regarding ACID transactions, what does \textbf{Isolation} mean?
\begin{choices}
  \choice It guarantees durability on disk.
  \choice It guarantees atomic rollback segments.
  \CorrectChoice Concurrent transactions are shielded according to the chosen isolation level so interactions are controlled.
  \choice It automatically shards data across nodes.
\end{choices}
\begin{solution}
  Isolation defines visibility/interaction rules (read committed, repeatable read, serializable, etc.) between overlapping transactions.
\end{solution}

\question[4]
Regarding ACID transactions, what does \textbf{Durability} mean?
\begin{choices}
  \choice It ensures logical schema validity after each statement.
  \CorrectChoice After a transaction commits, its effects survive subsequent failures (assuming storage honors the guarantee).
  \choice It ensures mutual exclusion of threads inside the JVM.
  \choice It implies infinite storage capacity.
\end{choices}
\begin{solution}
  Durability is the ``once committed, stays committed'' property modulo storage subsystem behavior.
\end{solution}

\question[4]
Spring Web is used primarily to handle arbitrary POJOs and register them as singleton services with no HTTP involvement.
\begin{choices}
  \choice TRUE
  \CorrectChoice FALSE
\end{choices}
\begin{solution}
  Spring Web centers on HTTP, servlets, MVC, WebClient/WebFlux, and related integrations---not on replacing the core container's general-purpose bean registration story.
\end{solution}

\question[4]
\textit{(Select all that apply.)} Which statements describe advantages of Spring's inversion of control?
\begin{checkboxes}
  \CorrectChoice It decouples components from concrete implementations by injecting collaborators.
  \choice It requires every class to manually \texttt{new} its dependencies for clarity.
  \CorrectChoice It simplifies testing by substituting mocks or stubs through the container or direct constructor injection.
  \choice It mandates tight coupling between layers to improve performance.
\end{checkboxes}
\begin{solution}
  IoC/DI promotes flexible wiring and test doubles; it discourages hard-coded \texttt{new} for every collaborator and reduces coupling rather than increasing it.
\end{solution}

\endgradingrange{csaspring}
